\centerline{\bf \Huge Домашнее задание.}
\centerline{\bf \Huge НОК, НОД, делимость и остатки}

\begin{flushright}
        \scriptsize{— Самое замечательное число\\
                    — 73. Вы скорее всего теряетесь в догадках почему?\\
                    — Нет!\\
                    — И в мыслях не было!\\
                    — 73\\
                    — Это двадцать первое простое число,\\ его зеркальное отражение 37 является двенадцатым,\\ чьё отражение 21 является результатом умножения — не упадите — семи и трёх!\\ Ну, не обманул?\\
                    Теория большего взрыва. Шелдон Купер}
\end{flushright}


\problem{1} Расул Владимирович загадал два числа и выписал на доску произведение этих чисел, равное $864$ и их НОД = $12$. Найдите НОК загаданных двух чисел?

\problem{2} Найдите хотя бы два числа А таких, что для них выполнены все три условия:\\
1) А при делении на $17$ дает остаток $1$\\
2) А при делении на $3$ дает остаток $0$\\
3) А отрицательное число\\

\problem{3} Начальник выписал на доску число, которое делится на $55$. Затем он решил стереть три цифры, и на доске осталось следующее число без трех цифр: $437$\_$70$\_$8$\_
Определите, какие цифры стер Начальник?
  
\problem{4} Известно, что число A при делении на $100$ дает остаток X, а число B остаток Y.\\
а) Делится ли A$-$B на $100$, если $X=Y+1$\\
б) Делится ли A+B на $100$, если $X+Y=9*10$\\
в) Делится ли A+B на $100$, если $X+Y=10*10$\\
г) Делится ли A$-$B на $100$, если $X=Y$\\